\section{Composite}
\label{sec:pat-composite}

\problemstatement{Compose objects into tree structures to represent
part-whole hierarchies, and let clients treat individual objects and
compositions of objects uniformly through the same interface.}

\begin{patternbox}[When You Reach for It]
The trigger is \emph{``a tree of objects where leaves and containers
should be handled the same way''} -- files and folders, a GUI of widgets
containing widgets, an org chart. When client code should not care whether
it holds one item or a group, use Composite.
\end{patternbox}

\subsection*{Understanding the pattern}

Define one \emph{component} interface implemented by both \emph{leaves}
(no children) and \emph{composites} (hold children, also components). A
composite's operation typically recurses over its children. Because both
share the interface, client code calls, say, \texttt{getSize()} on
anything and the recursion through the tree happens transparently.

\begin{ideabox}[Leaf and Container Share One Interface]
Make the container a component too, holding a list of components. An
operation on a container delegates to its children and aggregates. This
uniformity is the whole point: no \texttt{if (isFolder)} branching in
client code.
\end{ideabox}

\subsection*{C++ implementation}

\begin{lstlisting}
#include <bits/stdc++.h>
using namespace std;

struct FileSystemNode {                        // component
    virtual int size() const = 0;
    virtual void print(string indent = "") const = 0;
    virtual ~FileSystemNode() = default;
};

struct File : FileSystemNode {                 // leaf
    string name; int bytes;
    File(string n, int b) : name(move(n)), bytes(b) {}
    int size() const override { return bytes; }
    void print(string indent) const override {
        cout << indent << name << " (" << bytes << ")\n";
    }
};

struct Folder : FileSystemNode {               // composite
    string name;
    vector<unique_ptr<FileSystemNode>> children;
    Folder(string n) : name(move(n)) {}
    void add(unique_ptr<FileSystemNode> c) { children.push_back(move(c)); }
    int size() const override {
        int total = 0;
        for (auto& c : children) total += c->size();   // recurse
        return total;
    }
    void print(string indent) const override {
        cout << indent << name << "/\n";
        for (auto& c : children) c->print(indent + "  ");
    }
};

int main() {
    auto root = make_unique<Folder>("root");
    root->add(make_unique<File>("a.txt", 100));
    auto sub = make_unique<Folder>("sub");
    sub->add(make_unique<File>("b.txt", 250));
    root->add(move(sub));

    root->print();
    cout << "Total size: " << root->size() << "\n";   // treated uniformly
    return 0;
}
\end{lstlisting}

\begin{complexitybox}[Trade-offs]
\textbf{Pros.} Uniform treatment of leaves and groups; new node types slot
in easily; recursion over the tree is natural. \textbf{Cons.} The shared
interface can become overly general (leaf-only or container-only methods
awkwardly appear on both); type safety of ``only folders have children''
is relaxed.
\end{complexitybox}

\begin{takeawaybox}
Composite models part-whole trees behind a single component interface so
clients treat one and many identically, with container operations
recursing over children. Recognise it wherever a hierarchy of ``things and
groups of things'' must be processed uniformly.
\end{takeawaybox}
